% PCCSL 507 Machine Learning Lab - Experiment 2 (NumPy only)
\documentclass[11pt,twoside]{article}

\usepackage[paper=letterpaper,top=0.7in,bottom=0.63in,left=0.75in,right=0.75in,headheight=24pt,headsep=0.14in,footskip=0.38in]{geometry}

\usepackage[T1]{fontenc}
\usepackage{mathptmx} % Times clone
\usepackage{xcolor}
\definecolor{primary}{HTML}{1F3864}
\usepackage{fancyhdr}
\usepackage{array}
\usepackage{colortbl}
\usepackage{parskip}
\setlength{\parindent}{0pt}
\setlength{\parskip}{4pt}
\setlength{\abovedisplayskip}{4pt}
\setlength{\belowdisplayskip}{4pt}
\setlength{\abovedisplayshortskip}{2pt}
\setlength{\belowdisplayshortskip}{2pt}
\usepackage{graphicx}
\usepackage{amsmath}
\usepackage{listings}
\lstset{basicstyle=\footnotesize\ttfamily, breaklines=true, frame=single, showstringspaces=false}

\pagestyle{fancy}
\fancyhf{}
\fancyhead[C]{\textcolor{primary}{\textbf{\small PCCSL 507 Machine Learning Laboratory Record}}}
\renewcommand{\headrulewidth}{0.75pt}
\renewcommand{\headrule}{\hbox to\headwidth{\color{primary}\leaders\hrule height \headrulewidth\hfill}}
\fancyfoot[C]{\footnotesize DEPARTMENT OF CSE \hfill CAPE COLLEGE OF ENGINEERING ALAPPUZHA \hfill Page No. \thepage}
\renewcommand{\footrulewidth}{0pt}

\newcommand{\expheading}[1]{{\fontsize{14}{17}\selectfont\textbf{#1}\par}}
\newcommand{\fullrule}{\noindent\rule{\textwidth}{0.6pt}\par}

\begin{document}
\sloppy

% Sheet-local tightening so it fits one right-side page
{\setlength{\parskip}{1pt}\setlength{\abovedisplayskip}{2pt}\setlength{\belowdisplayskip}{2pt}
\begin{center}
\expheading{EXPERIMENT NO. 2}
\end{center}
\vspace{2pt}

\noindent\textbf{Date:} \underline{\hspace{2.6cm}} \hfill \textbf{Page No.:} \underline{\hspace{2.6cm}}\par
\vspace{6pt}

\expheading{TITLE}
\vspace{4pt}
\noindent NumPy Basics -- Student Marks Analysis\par
\vspace{4pt}

\expheading{OBJECTIVE}
\vspace{2pt}
\fullrule
\vspace{4pt}
\noindent To familiarize with NumPy array equations by storing the marks of 3 students in 5 subjects and finding the total, average, highest score and subject-wise average. \par
\vspace{4pt}

\expheading{AIM}
\vspace{4pt}
\noindent To store the marks of 3 students in 5 subjects in a NumPy array and compute the total, average, highest score and subject-wise average. \par
\vspace{4pt}

\expheading{THEORY}
\vspace{4pt}
Marks matrix ($3$ students $\times$ $5$ subjects):
\[M=\begin{bmatrix}m_{11} & \cdots & m_{15}\\m_{21} & \cdots & m_{25}\\m_{31} & \cdots & m_{35}\end{bmatrix}\]
Total of all marks:
\[T=\sum_{i=1}^{3}\sum_{j=1}^{5}m_{ij}\qquad\text{(np.sum)}\]
Overall average:
\[\bar{A}=\frac{T}{15}\qquad\text{(np.mean)}\]
Highest score and its position:
\[H=\max(M)\qquad\text{(np.max, np.argmax)}\]
Student-wise total (row sum, axis=1):
\[T_i=\sum_{j=1}^{5}m_{ij}\]
Subject-wise average (column mean, axis=0):
\[\bar{S}_j=\frac{1}{3}\sum_{i=1}^{3}m_{ij}\]
\vspace{4pt}

\expheading{PROCEDURE}
\vspace{4pt}
\noindent 1. Store the marks of 3 students in 5 subjects in a $3\times5$ NumPy array.\par
\noindent 2. Find the total (np.sum) and the average (np.mean) of all marks.\par
\noindent 3. Find the highest score (np.max) and its student/subject (np.argmax).\par
\noindent 4. Find the student-wise total (sum, axis=1) and subject-wise average (mean, axis=0).\par
\noindent 5. Display all results and verify them. \par
\vspace{4pt}

\expheading{DATASET DESCRIPTION}
\vspace{4pt}
\noindent Dataset Name: Student marks (synthetic, entered in program)\par
\noindent Source: 3 students $\times$ 5 subjects, marks out of 100\par
\noindent Number of Samples: 15 marks \hspace{0.3cm} Number of Features: 5 subjects\par
\noindent Target / Class: Not applicable (descriptive statistics)\par
\noindent Description: S1=[78, 85, 92, 88, 76]; S2=[65, 72, 80, 69, 74]; S3=[90, 88, 95, 91, 89]. \par
\vspace{4pt}
} % end sheet-local tightening

\newpage
\expheading{PROGRAM}
\vspace{6pt}

\begin{lstlisting}
import numpy as np


marks = np.array([
    [78, 85, 92, 88, 76],
    [65, 72, 80, 69, 74],
    [90, 88, 95, 91, 89],
])
print("Marks (3 students x 5 subjects):")
print(marks)

total = marks.sum()
average = marks.mean()
highest = marks.max()
pos = np.unravel_index(marks.argmax(), marks.shape)
print("Total =", total)
print("Average =", round(float(average), 2))
print("Highest score =", highest, "by student", pos[0] + 1,
      "in subject", pos[1] + 1)

stud_total = marks.sum(axis=1)
print("Student-wise total =", stud_total)
sub_avg = marks.mean(axis=0)
print("Subject-wise average =", np.round(sub_avg, 2))
\end{lstlisting}

\newpage
\expheading{OUTPUT}
\vspace{2pt}
\noindent\rule{\textwidth}{0.5pt}\par
\vspace{4pt}
\begin{verbatim}
Marks (3 students x 5 subjects):
[[78 85 92 88 76]
 [65 72 80 69 74]
 [90 88 95 91 89]]
Total = 1232
Average = 82.13
Highest score = 95 by student 3 in subject 3
Student-wise total = [419 360 453]
Subject-wise average = [77.67 81.67 89.   82.67 79.67]
\end{verbatim}
\vspace{4pt}

\expheading{RESULT}
\vspace{4pt}
{\fontsize{12}{15}\selectfont
The above program was implemented and executed successfully using Python, and the corresponding output was obtained and verified. The results of the \textbf{NumPy array statistics} were analyzed and found to be satisfactory, thereby achieving the desired objective of the experiment.\par
}

\vspace{12pt}

\noindent\textbf{Faculty Signature:} \underline{\hspace{3.6cm}} \hfill \textbf{Date:} \underline{\hspace{2.8cm}}\par

\end{document}
